\documentclass[11pt]{article}

\usepackage[left=0.6 in,top=0.4in,right=0.6in,bottom=0.4in]{geometry}
\usepackage[parfill]{parskip} % Remove paragraph indentation

\usepackage{enumitem}
\usepackage{hyperref}
\usepackage{amsmath, amsfonts}

\begin{document}
\thispagestyle{empty}

\begin{center}
  \bf \Huge  First Last, Ph.D.
  \vspace{.1em}
  \hrule
\end{center}

\begin{center}
  \vspace{.4em}
  702-354-0866 $\Big \Vert$
  \href{mailto:firstlast@gmail.com}{firstlast@gmail.com} $\Big \Vert$ 
  \href{https://www.linkedin.com/in/first-last/}{linkedin.com/in/first-last} $\Big \Vert$
  \href{https://github.com/firstlast}{github.com/firstlast}
\end{center}

\textbf{\Large Experience}
\vspace{.3em}
\hrule
\vspace{.3em}

\textbf{\large Senior Data Scientist} \hfill \textit{June 2022 - Present}\\
\textit{Company name} 

\begin{itemize} \setlength\itemsep{-0.5em}
  \item Developed a YOLO-based object-detection model in PyTorch that achieved a 
    80\% mAP evaluation on a in-house thermal-image vehicle dataset. 

  \item Designed an anomaly-detection model in PyTorch using an LSTM 
    variational-autoencoder to detect component failure in a pulse-power x-ray 
    machine. The model accurately predicted the location of a component failure with 
    87\% accuracy and cut machine repair times by 50\%.

  \item Applied the U-Net architecture with PyTorch to develop a deblurring model for 
    a pulse-powered x-ray machine. The model was applied by the machine operators 
    and reduced machine configuration time by 30\%.

  \item Facilitated the scaling of locally developed models to be trained on our internal
    high-performance computing (HPC). Would take code written by the other 
    Data Scientists and Analysts and would refactor so that they could run in a 
    multi node, multi gpu setting. 

  \item Developed interactive dashboards using Streamlit to give non-technical staff the
    abilty to interact with the capabilities designed by our team of Data Scientists.

  \item Mentored a summer intern in the development and execution of models on a 
    HPC cluster, involving multi-node, multi-GPU setups and SLURM resource management.

  \item Sole developer of a Django web interface, which wrapped Neo4j. The interface
    gave non-technical staff the ability to query and update the company's Neo4j
    graph-database without needing to know any query syntax.

\end{itemize}


\vspace{.3em}
\textbf{\Large Technical Skills}
\vspace{.3em}
\hrule
\vspace{.3em}

\textbf{Programming Languages:}
Python, SQL, Bash, Cypher, Lua, \LaTeX

\textbf{Libraries:} 
PyTorch, Scikit-Learn, Pandas, Scipy, Sqlalchemy, Matplotlib, Plotly, Streamlit, Django
 

\textbf{Developer Tools:} Git, AWS, Unix/Linux, Slurm, Vim, Tmux, Neo4j


\vspace{.3em}
\textbf{\Large Education}
\vspace{.3em}
\hrule
\vspace{.3em}

\textbf{University 2 - Ph.D.} \hfill \textit{June 2022} \\
Research Area: Probability (Stochastic Processes and Partial Differential Equations) \\
Advisor: \href{http://}{First Last}

\textbf{University 1 - B.S./M.S.} 
\hfill \textit{December 2015 / January 2018} \\
Major: Applied Mathematics

\vspace{.3em}
\textbf{\Large Publications}
\vspace{.3em}
\hrule
\vspace{.3em}

\begin{itemize}[leftmargin=*]
  \item[1:] Q1 probability and stochastic process journal

  \item[2:] Q2 probability journal

\end{itemize}

\end{document}
